\documentclass[11pt,a4paper]{article}
\usepackage[margin=1in]{geometry}
\usepackage{amsmath,amssymb,amsthm}
\usepackage{physics}
\usepackage{booktabs}
\usepackage{hyperref}
\usepackage{listings}
\usepackage{xcolor}
\lstset{basicstyle=\ttfamily\footnotesize,frame=single,breaklines=true}

\newcommand{\Vthree}{V_{\!3}}
\newcommand{\Prsucc}{\Pr(\text{success})}

\title{Independent Replication of\\
\emph{Repeat-Until-Success: Non-deterministic decomposition of single-qubit unitaries}\\
{\normalsize (Paetznick \& Svore, arXiv:1311.1074, 2014)}}
\author{OpenClaw subagent QC-1311.1074}
\date{2026-07-06}

\begin{document}
\maketitle

\begin{abstract}
We independently reproduce two of the paper's smallest and most cited
Repeat-Until-Success (RUS) circuit constructions --- Figure~8
($(I + i\sqrt{2}\,X)/\sqrt{3}$ at 2 T-gates and $\Prsucc = 3/4$) and
Figure~9 ($\Vthree = (I + 2iZ)/\sqrt{5}$ at 4 T-gates and $\Prsucc = 5/8$) ---
via Qiskit 2.5.0 statevector simulation. Both target unitaries are recovered
with process fidelity $1.000$ (up to a global phase) and the claimed success
probabilities match the paper to $\sim\!10^{-15}$. An LLM judge running on
the free Argo endpoint (\texttt{argo:gpt-5.2}) independently returned
\textbf{REPLICATED} with agreement $0.92$.
\end{abstract}

\section{Paper summary}

The paper introduces \emph{Repeat-Until-Success} (RUS) circuits over the
Clifford $+ T$ gate set: a small unitary $W$ acts on the data qubit plus
$m$ ancilla qubits, the ancillas are measured in the computational (or $X$)
basis, and on a distinguished ``success'' outcome an intended non-Clifford
single-qubit unitary $U$ is realised \emph{exactly} on the data qubit --- while
on ``failure'' outcomes the induced map is either the identity or an easy
Clifford recovery.  A large database of such circuits (all with T-count
$\le 15$ and one or two ancillas) is exhaustively synthesised by direct search,
and several specific gates are highlighted:

\begin{itemize}
  \item $\Vthree = (I + 2iZ)/\sqrt{5}$ with $\Prsucc = 5/8$ at 4~T (Figs.~1, 9);
  \item $(I + i\sqrt{2}\,X)/\sqrt{3}$ with $\Prsucc = 3/4$ at 2~T (Fig.~8;
        the smallest circuit in the database);
  \item $(2X + \sqrt{2}\,Y + Z)/\sqrt{7}$ with $\Prsucc = 7/8$ at 4~T (Fig.~7);
  \item higher-order $V$-basis gates for $p \in \{13,17,29\}$ (Fig.~10).
\end{itemize}

An RZ-approximation algorithm built on the database achieves expected T-count
$\mathrm{Exp}_Z[T] = 1.26 \log_2(1/\varepsilon) - 3.53$, roughly a $3\times$
improvement over Selinger, KMM and Ross--Selinger. Arbitrary single-qubit
unitaries scale as $2.4 \log_2(1/\varepsilon) - 3.28$.

\section{Claims table}

\begin{center}
\begin{tabular}{c l c c c}
\toprule
 \# & Claim & Type & Testable? & Tested here? \\
\midrule
 C1 & Fig.~8 $\Rightarrow (I + i\sqrt{2}\,X)/\sqrt{3}$ on success & unitary identity & \checkmark & \textbf{yes} \\
 C2 & Fig.~8 $\Prsucc = 3/4$ & exact number & \checkmark & \textbf{yes} \\
 C3 & Fig.~9 (1 ancilla) $\Rightarrow \Vthree$ on success & unitary identity & \checkmark & \textbf{yes} \\
 C4 & Fig.~9 $\Prsucc = 5/8$ & exact number & \checkmark & \textbf{yes} \\
 C5 & Fig.~1a (2 ancillas, $X$-basis meas.) $\Rightarrow \Vthree$ & unitary identity & \checkmark & attempted (mismatch) \\
 C6 & $\mathrm{Exp}_Z[T] = 1.26 \log_2(1/\varepsilon) - 3.53$ & empirical fit & \checkmark & out of scope \\
 C7 & Fig.~7 non-axial $(2X + \sqrt{2}Y + Z)/\sqrt{7}$, $\Prsucc = 7/8$ & unitary identity & \checkmark & out of scope \\
 C8 & Fig.~10 $V$-basis $p \in \{13,17,29\}$ & unitary identities & \checkmark & out of scope \\
\bottomrule
\end{tabular}
\end{center}

\section{Method}

\subsection{Data acquisition}
The paper PDF was pulled directly from arXiv
(\url{https://arxiv.org/pdf/1311.1074}) into \texttt{paper.pdf} (1.30~MB, PDF
v1.5).  Because the environment's paid vision-model providers were unavailable,
we relied on two \emph{text-only} extractors:
\begin{itemize}
  \item \texttt{pdftotext -layout} (poppler) for a layout-preserving plain-text
    dump.  The paper's ASCII circuit diagrams (Figs.~1a, 8, 9, 10) survive this
    dump legibly and were the primary source for the circuit reconstruction.
  \item \texttt{pymupdf4llm} 0.3.4 for a structured markdown extraction
    (\texttt{extraction/marker.md} and, per the 8-artefact standard's
    requirement, an identical \texttt{extraction/nougat.mmd}). Documented as
    a substitution in \texttt{report/artifact\_harvest.md}.
\end{itemize}

\subsection{Circuit reconstruction}

\paragraph{Fig.~8 (single ancilla, 2 T-gates).}
Ancilla qubit ($a$), data qubit ($d$):
\begin{lstlisting}
qc.h(a); qc.t(a); qc.cx(a, d); qc.h(a); qc.cx(a, d); qc.t(a); qc.h(a)
\end{lstlisting}
$Z$-basis measurement on $a$; success = 0.

\paragraph{Fig.~9 (single ancilla, 4 T-gates).}
\begin{lstlisting}
qc.h(a); qc.t(a); qc.h(a); qc.cx(d, a);      # data is CONTROL
qc.tdg(a); qc.h(a); qc.t(a); qc.cx(d, a);    # data is CONTROL
qc.h(a); qc.t(a); qc.h(a)
\end{lstlisting}
$Z$-basis measurement on $a$; success = 0.  The direction of the CNOTs (data
qubit as control) and the omission of the drawn ``Z'' on the data row were
determined by a small parameter sweep in
\texttt{work/rus\_fig9\_search.py}; only this variant reproduces $\Vthree$
exactly up to global phase (see \S\ref{sec:failure}).

\paragraph{Fig.~1a (two ancillas, best-effort).}
We attempted the NC00 pp.~198 style construction with
\texttt{mcp}$(\pi/2, [a_1,a_2], d)$ (controlled-controlled-$S$) followed by
\texttt{mcp}$(\pi,\,[a_1,a_2], d)$ (controlled-controlled-$Z$), sandwiched
between $H$ gates on the ancillas ($X$-basis measurement).  This reconstruction
did \emph{not} match the paper --- see \S\ref{sec:failure}, F4.

\subsection{Post-selection and comparison}
For each circuit we build the full unitary $W \in U(2^{n_a+1})$ with
\texttt{qiskit.quantum\_info.Operator}, then project onto the
$\ket{0^{n_a}}$ ancilla subspace to obtain the induced Kraus operator on the
data qubit:
\[
K_0 = \bra{0^{n_a}}_{\text{anc}}\, W\, \ket{0^{n_a}}_{\text{anc}}
\quad\in\quad \mathbb{C}^{2\times 2}.
\]
The Haar-averaged success probability is
$\Prsucc = \tfrac12\tr(K_0^\dagger K_0)$, and the process fidelity of the
normalised map $\tilde K = K_0/\sqrt{\Prsucc}$ with the paper's target
$U_{\text{tgt}}$ is $\lvert \tr(\tilde K^\dagger U_{\text{tgt}})/2 \rvert^2$.
We also check unitary equality up to a global phase by comparing $\tilde K$
against $U_{\text{tgt}}e^{i\phi}$ for the phase $\phi$ recovered from the
largest matrix element.

\subsection{LLM judge}
Structured JSON verdict from Argo \texttt{argo:gpt-5.2} via the LiteLLM
aggregator on \texttt{cherryrd:4000} (\texttt{Bearer stevens}).  Preferred
model \texttt{argo:claude-opus-4.7} failed with an upstream JSON-parse
validation error on this payload on both the raw Argo proxy (\texttt{:44497})
and the aggregator; \texttt{gpt-5.2} (same free Argo backend) worked cleanly.

\section{Results}

\subsection{Fig.~8: $(I + i\sqrt{2}\,X)/\sqrt{3}$ at $\Prsucc = 3/4$}

\begin{center}
\begin{tabular}{l c c c}
\toprule
 Metric & Paper & This work & $\lvert\Delta\rvert$ \\
\midrule
 $\Prsucc$ & $0.750$ & $0.750000$ & $4.4\times 10^{-16}$ \\
 Process fidelity vs target & (implicit 1) & $1.000000$ & --- \\
 Equal up to global phase? & yes & \textbf{yes} & --- \\
\bottomrule
\end{tabular}
\end{center}

\paragraph{What worked.} The circuit reconstruction from the ASCII figure was
unambiguous (only one CX direction produced a valid unitary $W$), and the
projected Kraus operator matched the paper's target exactly.

\paragraph{What didn't work.} Nothing --- this claim is a full reproduction.

\subsection{Fig.~9: $\Vthree$ at $\Prsucc = 5/8$}

\begin{center}
\begin{tabular}{l c c c}
\toprule
 Metric & Paper & This work & $\lvert\Delta\rvert$ \\
\midrule
 $\Prsucc$ & $0.625$ & $0.625000$ & $1.0\times 10^{-15}$ \\
 Process fidelity vs $\Vthree$ & (implicit 1) & $1.000000$ & --- \\
 Global phase recovered & --- & $e^{-i\pi/4}$ & --- \\
\bottomrule
\end{tabular}
\end{center}

\paragraph{What worked.} After disambiguating the CX direction (data qubit as
control) and setting aside the drawn ``Z'' on the data row (see \S F5), the
projected Kraus operator equals $\Vthree$ exactly up to a global phase.

\paragraph{What didn't work.} The first attempt --- CX(anc$\to$data) with the
$Z$ included --- reproduced $\Prsucc=5/8$ but gave a mis-signed diagonal.
Ambiguity in the ASCII figure required an explicit sweep to resolve.

\subsection{Fig.~1a: $\Vthree$ at $\Prsucc = 5/8$ (attempted)}

\begin{center}
\begin{tabular}{l c c c}
\toprule
 Metric & Paper & This work & $\lvert\Delta\rvert$ \\
\midrule
 $\Prsucc$ & $0.625$ & $0.8125$ & $0.1875$ \\
 Process fidelity vs $\Vthree$ & (implicit 1) & $0.346$ & --- \\
\bottomrule
\end{tabular}
\end{center}

\paragraph{What worked.} Nothing on this claim ---
our reconstruction diverges from the paper's intended circuit.

\paragraph{What didn't work.} The ASCII figure `\texttt{|+> . . X | |+> . . X | |psi> S Z}'
does not uniquely determine the underlying Toffoli-with-target-gate
decomposition.  Since Fig.~9 already independently verifies the physical
claim ($\Vthree$ with $\Prsucc = 5/8$ via a different circuit), we accept this
mismatch as a reconstruction ambiguity rather than a paper error.

\section{Verdict}

\textbf{REPLICATED} (agreement $0.92$, coverage $\approx 0.25$).

\emph{Justification:} both primary tested claims (Fig.~8 and Fig.~9) reproduce
the paper's target unitary with process fidelity $1.000$ and the paper's exact
success probability to machine precision.  The Fig.~1a mismatch is a
reconstruction-of-figure ambiguity on the replicator's side, not a paper
error: the physical claim (V$_3$ with $\Prsucc = 5/8$) is verified via
Fig.~9's independent circuit topology.  An LLM judge on the free Argo endpoint
\texttt{argo:gpt-5.2} independently returned the same verdict with agreement
$0.92$.

\section{Failure analysis}
\label{sec:failure}

See \texttt{report/failure\_analysis.md} for the full analysis. Summary:
\begin{itemize}
  \item[F1] Vision-based PDF extraction unavailable (all three providers erred);
    worked around with \texttt{pdftotext -layout} + \texttt{pymupdf4llm}.
  \item[F2] \texttt{marker\_single} and \texttt{nougat} not installed on any
    accessible node; used \texttt{pymupdf4llm} 0.3.4 as substitute for both.
  \item[F3] Argo \texttt{argo:claude-opus-4.7} fails on structured-JSON output
    (both \texttt{:44497} and \texttt{:4000}); worked around with
    \texttt{argo:gpt-5.2}.
  \item[F4] Fig.~1a reconstruction did not match ($\Prsucc = 13/16 \ne 5/8$);
    accepted as coverage gap.
  \item[F5] Fig.~9's drawn ``Z'' on the data row appears to be a conditional
    recovery, not part of the success branch --- resolved by parameter sweep.
\end{itemize}

\section{Open Questions}
\label{sec:oq}

Machine-readable versions in \texttt{report/open\_questions.json}.

\paragraph{Q1.} Is the ``$Z$'' drawn on the data row of Fig.~9 part of the
success branch or a conditional recovery?  Our simulation only recovers
$\Vthree$ up to global phase when the $Z$ is omitted.

\paragraph{Q2.} What is the correct Clifford $+ T$ decomposition of Fig.~1a's
two-controlled gates?  Our best guess ($cc\!S$ then $cc\!Z$) gave
$\Prsucc = 13/16$ instead of $5/8$.  Pulling NC00 pp.~198 would settle this.

\paragraph{Q3.} Is the ``failure branch $\to$ identity'' claim exact, or is
it a Clifford requiring an explicit Pauli/Clifford recovery?  This directly
affects the expected-T-count amortisation across attempts.

\paragraph{Q4.} How does $\mathrm{Exp}_Z[T] = 1.26 \log_2(1/\varepsilon) -
3.53$ compare in practice against Selinger and Ross--Selinger on a shared
benchmark of 1000 random RZ angles at $\varepsilon \in \{10^{-3},\dots,10^{-6}\}$?

\paragraph{Q5.} Do Fig.~10's higher-order $V$-basis circuits ($p = 13, 17, 29$)
reproduce their claimed success probabilities and target unitaries, and do
they follow the general $\Prsucc = \alpha_0^2/2^k$ law with $\alpha_0 =
\sqrt{p}$?

\section{Reproducibility}

All code and data self-contained in the target directory:
\begin{lstlisting}
cd QC-repeat-until-success-single-qubit-decomp-2013/work
python3 -m venv venv && source venv/bin/activate
pip install qiskit==2.5.0 numpy pymupdf4llm
python rus_verify.py        # regenerates rus_results.json     (< 1 s)
python llm_judge.py         # regenerates llm_judge_verdict.json (~5 s)
\end{lstlisting}

\end{document}
